\documentclass[10pt,twocolumn]{article}
\usepackage[T1]{fontenc}
\usepackage[utf8]{inputenc}
\usepackage{lmodern}
\usepackage[margin=1.8cm,top=2.4cm,bottom=2cm,columnsep=0.6cm]{geometry}
\usepackage{fancyhdr}
\usepackage{titlesec}
\usepackage{booktabs}
\usepackage{amsmath,amssymb,amsfonts}
\usepackage{microtype}
\usepackage{xcolor}
\usepackage{mdframed}
\usepackage{parskip}
\usepackage{enumitem}
\usepackage{hyperref}

\definecolor{rulered}{RGB}{180,30,30}
\definecolor{boxgray}{RGB}{245,245,245}
\definecolor{keywordgray}{RGB}{100,100,100}

\pagestyle{fancy}
\fancyhf{}
\fancyhead[L]{\textsc{\small Claude Mag}}
\fancyhead[C]{\small\textit{Machine Learning Research Digest}}
\fancyhead[R]{\small 2026-10-28}
\fancyfoot[C]{\small\thepage}
\renewcommand{\headrulewidth}{0.8pt}

\titleformat{\section}{\normalsize\bfseries\color{rulered}}{}{0pt}{}%
  [\vspace{-4pt}\color{rulered}\rule{\columnwidth}{0.4pt}]
\titlespacing{\section}{0pt}{8pt}{4pt}

\hypersetup{colorlinks=true,urlcolor=rulered,linkcolor=black}

\begin{document}
\setlength{\parindent}{0pt}
\setlength{\parskip}{4pt}

{\small\color{keywordgray}\textbf{Keywords:}
  multi-objective reinforcement learning \textbullet{}
  ensemble policies \textbullet{}
  trajectory attribution \textbullet{}
  objective conflict}\par\vspace{4pt}

{\LARGE\bfseries\raggedright
  Multi-Objective Reinforcement Learning Using Routed
  Ensembles and Trajectory Attribution}\par\vspace{4pt}

{\large\itshape\raggedright
  Objective-Conflict Coefficients Quantify Preference Trade-offs, While
  Routed Ensembles and Kernel-Surrogate Attribution Close the
  Multi-Objective Efficiency Gap}\par\vspace{6pt}

{\small\textbf{By} Zhenshuo Zhang, Ruoxuan Xiong, Hongyang R.\ Zhang}\par\vspace{2pt}

{\small\color{keywordgray}arXiv:2610.35733 \textbullet{}
  NeurIPS 2026}
\par\vspace{2pt}\noindent{\color{rulered}\rule{\columnwidth}{0.6pt}}\vspace{6pt}

Multi-objective reinforcement learning demands a single agent to
simultaneously satisfy multiple reward objectives weighted by a
user-supplied preference vector---a task made hard when objectives
conflict and gradient directions oppose each other. This paper
introduces objective-conflict coefficients that formally quantify
this conflict, derives matching online regret bounds, and proposes
two practical algorithms: a routed ensemble that routes each trajectory
to the basis policy with highest expected return under its preference
weight, and a trajectory attribution module based on kernel ridge
regression that identifies which trajectories are most valuable for
which preference objectives.

\begin{mdframed}[backgroundcolor=boxgray,linecolor=rulered,linewidth=1pt,
    innerleftmargin=6pt,innerrightmargin=6pt,
    innertopmargin=6pt,innerbottommargin=6pt]
\textbf{\small AT A GLANCE}\par\vspace{4pt}
\begin{description}[leftmargin=0pt,labelsep=4pt,itemsep=2pt,parsep=0pt,
    font=\small\bfseries]
  \item[Problem:] A single preference-conditioned policy fails when
    objectives strongly conflict; the degree of conflict was not
    previously formalised.
  \item[Why it matters:] MORL is central to robotics, resource
    allocation, and fairness-aware learning, where trade-offs among
    objectives must be navigated at deployment.
  \item[Method:] Objective-conflict coefficients + routed ensemble
    policy learning (trajectory routing to K basis policies) +
    kernel ridge regression trajectory attribution.
  \item[Result:] Consistent improvements over single-policy and
    standard ensemble baselines on multi-objective robotic control
    benchmarks.
  \item[Evidence:] Multi-objective MuJoCo locomotion benchmarks;
    formal regret analysis as a function of maximum conflict
    coefficient.
\end{description}
\end{mdframed}

\section{The Big Picture}

In multi-objective reinforcement learning (MORL), the agent receives a
vector reward $\mathbf{r} = (r_1, \ldots, r_K)$ and a preference
weight $\mathbf{w} \in \Delta^{K-1}$. The goal is to find a policy
$\pi$ that maximises $\mathbf{w}^\top J(\pi)$ across many possible
$\mathbf{w}$, effectively learning the entire Pareto front of
preference-weighted returns.

Conflict arises when $\nabla_\theta J_i(\pi_\theta)$ and
$\nabla_\theta J_j(\pi_\theta)$ point in opposing directions, forcing
any gradient step to damage at least one objective. Prior approaches
either concatenate $\mathbf{w}$ to the state (ignoring conflict) or
use hand-crafted ensemble partitions (ignoring data efficiency). This
work addresses both gaps.

\section{Why Existing Approaches Fall Short}

\textbf{Conditioned single policy:}
Appending $\mathbf{w}$ to the state does not prevent gradient conflicts.
For strongly conflicting objectives, the single policy converges to a
suboptimal compromise across all preferences rather than achieving the
Pareto front.

\textbf{Na\"ive ensembles:}
Manually partitioning the preference space into $K$ fixed regions assigns
trajectories to basis policies by region membership, ignoring that
trajectories generated under one preference weight may be highly
informative for another.

\textbf{No formal measure of conflict:}
Without a quantitative measure of objective conflict, there is no
principled way to predict how many basis policies are needed or which
trajectory assignments are most beneficial.

\section{Core Mechanism}

\textbf{Routed Ensemble Policy Learning} maintains $K$ basis policies
$\{\pi_1, \ldots, \pi_K\}$ and iterates:
\begin{enumerate}[itemsep=2pt,parsep=0pt]
  \item Collect trajectories $\{\tau_j\}$ by sampling basis policies
    according to the current preference weight distribution.
  \item \textbf{Route:} Assign each $\tau_j$ to the basis policy with
    the highest expected return under $\tau_j$'s preference weight:
    $k^*(\tau_j) = \arg\max_k \mathbf{w}_j^\top J_k(\pi_k)$.
  \item \textbf{Update:} Apply PPO to each $\pi_k$ using only
    trajectories routed to it.
\end{enumerate}

\textbf{Trajectory Attribution} uses a kernel ridge regression surrogate
$\hat{f}(\pi, \mathbf{w}; \mathcal{D})$ to estimate policy performance
as a function of the training trajectory set $\mathcal{D}$.
Leave-one-out attribution scores $\phi_j^k$ identify which trajectories
help or hurt each basis policy under each preference weight, enabling
curriculum sampling.

\section{Technical Formulation}

\textbf{Objective-Conflict Coefficient.} For objectives $i$ and $j$:
\[
c_{ij}(\pi_\theta) = \frac{-\nabla_\theta J_i^\top \nabla_\theta J_j}
  {\|\nabla_\theta J_i\|\,\|\nabla_\theta J_j\|}
\]
A positive $c_{ij}$ indicates conflict (gradients oppose); the maximum
over all pairs and all preference weights bounds the MORL regret:
\[
\text{Regret}(T) = O\!\left(\sqrt{T \cdot c_{\max}}\right)
\]
where $c_{\max} = \max_{i \neq j, \mathbf{w}} c_{ij}(\pi^*(\mathbf{w}))$.

\textbf{Kernel Surrogate.} Let $\phi(\tau, \mathbf{w})$ be a feature
embedding of trajectory $\tau$ and preference weight $\mathbf{w}$.
The surrogate is:
\[
\hat{f}(\mathbf{w}; \mathcal{D}) =
  \mathbf{y}^\top (K + \lambda I)^{-1} k(\mathbf{w})
\]
where $K_{ij} = \kappa(\phi(\tau_i), \phi(\tau_j))$,
$\mathbf{y}_j = \mathbf{w}_j^\top J(\pi)$, and $\lambda$ is a
regularisation constant. The attribution of trajectory $\tau_j$ for
preference $\mathbf{w}$ is the leave-one-out change in $\hat{f}$.

\section{Learning or Inference Procedure}

Each training step:
\begin{enumerate}[itemsep=2pt,parsep=0pt]
  \item Sample a batch of preference weights $\{\mathbf{w}_j\}_{j=1}^B$
    from the prior distribution over $\Delta^{K-1}$.
  \item Collect trajectories by rolling out each $\pi_k$ under its
    currently assigned preference weights.
  \item Compute attribution scores; upsample high-attribution trajectories.
  \item Route trajectories to basis policies and apply PPO updates.
\end{enumerate}

At test time, given a user preference $\mathbf{w}^*$, the system
evaluates $J_k(\pi_k)$ for each basis policy and returns the best.

\section{What the Guarantee Says}

The regret bound $O(\sqrt{T \cdot c_{\max}})$ provides an
instance-dependent characterisation: MORL is hard when objectives
conflict strongly ($c_{\max}$ large) and easy when they are compatible
($c_{\max}$ small). The routed ensemble achieves this rate for any
$K \geq c_{\max} / \epsilon$ for tolerance $\epsilon$.
This is the first formal characterisation of MORL hardness in terms
of objective geometry.

\section{Experimental Findings}

\begin{itemize}[itemsep=2pt,parsep=0pt]
  \item \textbf{Multi-objective MuJoCo locomotion:} Routed ensemble
    outperforms conditioned single-policy baselines by $8$--$15\%$
    in hypervolume metric across 4 multi-objective task suites.
  \item \textbf{Versus na\"ive ensemble:} Attribution-based curriculum
    sampling improves sample efficiency by $\sim$30\% (same
    hypervolume with fewer environment interactions).
  \item \textbf{Conflict coefficient:} Empirical $c_{\max}$ correlates
    strongly with the observed performance gap between single-policy
    and ensemble methods ($r = 0.81$), validating the theoretical
    predictor.
  \item \textbf{Scalability:} Performance improvements hold for
    $K \in \{2, 4, 8\}$ basis policies; gains plateau at $K = 8$.
\end{itemize}

\section{Ablations and Interpretation}

\textbf{Routing vs.\ random assignment:}
Random assignment of trajectories to basis policies reduces hypervolume
by 7\%, confirming that conflict-aware routing is the key driver.

\textbf{Attribution vs.\ uniform sampling:}
Replacing kernel attribution with uniform trajectory sampling costs
$\sim$20\% in sample efficiency, isolating the contribution of the
surrogate.

\textbf{Conflict coefficient threshold:}
Task suites with $c_{\max} < 0.1$ show no significant gain from ensembles;
the single-policy baseline already matches. This provides a practical
criterion for when ensembles are needed.

\section*{Reference}

Zhenshuo Zhang, Ruoxuan Xiong, Hongyang R.\ Zhang.
\textbf{Multi-Objective Reinforcement Learning Using Routed Ensembles
and Trajectory Attribution}.
arXiv:2610.35733. NeurIPS 2026.
\url{https://arxiv.org/abs/2610.35733}

\end{document}
